% !PRINTED 725
\weberpage{725}
%<<CONT
and consequently, by (10),
\[
N\varphi(\Theta) = 2^{\nu'-1}\,N\psi(\Theta).
\]
From this there results
\weq{13}{A = N\psi(\Theta),}
and from this it follows, having regard to (12), that $A$ is an \gesperrt{odd
integer.}

If we denote the number $A$ belonging to $m = 2^\varkappa$ by $A_\varkappa$, then
from \S.\,197, (4) there results, by complete induction, the
expression for the first factor of the class number:
\weq{14}{k = A_4\,A_5\,\dots\,A_\varkappa,}
and from this the theorem:

1. \gesperrt{The first factor of the class number is an odd integer.}

The number $A_\varkappa$ is, by formula (10), not hard to compute for the
first values of $\varkappa$.

For $m = 16$ one finds $\nu' = 2$, $\beta = 0, 1$, $c_0 = c_1 = 1$
consequently
\[
\psi(\Theta) = 1, \quad A_4 = 1.
\]

For $m = 32$, $\nu' = 4$, $\beta = 0, 1, 2, 3$.
\[
c_0 = 1, \quad c_1 = 1, \quad c_2 = -1, \quad c_3 = -1,
\]
\[
\psi(\Theta) = -\Theta^2,
\]
hence also $A_5 = 1$.

For $m = 64$, $\nu' = 8$, $\beta = 0, 1, 2, 3, 4, 5, 6, 7$ the
absolutely least residues of $5^\beta$ are
\[
1,\ 5,\ 25,\ -3,\ -15,\ -11,\ 9,\ -19;
\]
so the $c_\beta$ are:
\[
+1, +1, +1, -1, -1, -1, +1, -1
\]
and
\[
\psi(\Theta) = -\Theta^3 + \Theta^6 - \Theta^7, \quad \psi(\Theta)\,\psi(-\Theta) = \Theta^6(\Theta^6 - 2i),
\]
from which it easily results:
\[
A_6 = 17.
\]
A somewhat longer calculation gives in addition
\[
A_7 = 21121.
\]
